\scanpage{183}\mspage{38}%
% archive sheet number: 183
% One continuous running-prose block; no visible paragraph indent or extra gap at Une fois.
dans cette première approche pré-catégorie dérivée, par la connaissance du groupe fondamental $\pi_1$, donc des catégories des coeff. \unsure{constant|constants} tordus $=\pi_1$-\unsure{modules|module}, les foncteurs
% L5: the argument separator has a high point (about 344,446) and a separate lower comma-shaped stroke.
% Preferred semicolon; alternative comma with an unassigned detached point. Neither element is omitted.
$H^i(-\unsure{;|,}M)$ sur ceux-ci avec \unsure{le|la} structure de \unsure{foncteur|foncteurs} \deleted{\unsure{ab|ab.}, de \unsure{foncteur|foncteurs}} \inserted{\unsure{cohomologique|cohomologiques}} relatif aux suites exactes, \unsure{le|la} \deleted{\unsure{struc}}structure de cup-produit etc --\raisebox{0pt}[0pt][0pt]{\parbox{0pt}{\figure{\begin{tikzpicture}[overlay,x=1em,y=1em]\fill (-0.55,-0.60) circle[radius=0.045em];\end{tikzpicture}}}} reliés par certaines propriétés formelles. Une fois qu'on dispose du langage des cat. dérivées, la \inserted{sous catégorie de la catégorie dérivée des} complexes abéliens, \unsure{formee|formée|form-} \inserted{aux complexes de} faisceaux de coh. loc constante sur $X$, avec sa structure triangulé et sa structure \deleted{\unsure{m|de|sou}}multiplicative
% L13: detached compact point below tensor, approximately (1165,1373); origin and function unassigned.
$\overset{\mathbb{L}}{\otimes}\raisebox{0pt}[0pt][0pt]{\parbox{0pt}{\figure{\begin{tikzpicture}[overlay,x=1em,y=1em]\fill (-0.45,-0.95) ellipse[x radius=0.05em,y radius=0.07em];\end{tikzpicture}}}}$ (et éventuellement $\mathbb{R}\mathrm{Hom}$ \ldots) \unsure{donnant|donnait} un candidat plus satisfaisant pour espérer récupérer le type d'homotopie. Je n'ai guère d'idée
% L17: separate short stroke below si (about 241,1836); two compact spots beside ils/permettent
% (about 802,1786 and 874,1810) have unassigned origin, irrespective of their colour.
si\raisebox{0pt}[0pt][0pt]{\parbox{0pt}{\figure{\begin{tikzpicture}[overlay,x=1em,y=1em]\draw[line width=0.6pt] (-0.45,-0.80) -- (-0.40,-0.97);\end{tikzpicture}}}} oui ou non ils\raisebox{0pt}[0pt][0pt]{\parbox{0pt}{\figure{\begin{tikzpicture}[overlay,x=1em,y=1em]\fill (0.08,0.58) circle[radius=0.04em];\end{tikzpicture}}}} \raisebox{0pt}[0pt][0pt]{\parbox{0pt}{\figure{\begin{tikzpicture}[overlay,x=1em,y=1em]\fill (-0.12,0.05) ellipse[x radius=0.055em,y radius=0.07em];\end{tikzpicture}}}}permettent de le récupérer \inserted{en effet}, mais par \raisebox{0pt}[0pt][0pt]{\parbox{0pt}{\figure{\begin{tikzpicture}[overlay,x=1em,y=1em]\fill (-0.10,-0.92) circle[radius=0.045em];\fill (0.06,-0.67) ellipse[x radius=0.045em,y radius=0.06em];\end{tikzpicture}}}}contre n'ai guère de doute qu'en allant au bout de la ``linéarisation'', c'est à dire en partant dans le
% L21 underline: first physical stroke about (318,2273)--(556,2261), second about (581,2262)--(673,2264).
% The gap lies mainly under l; a short overlay preserves the second stroke's beginning under l's terminal extent.
\uline{non abé}l\raisebox{0pt}[0pt][0pt]{\parbox{0pt}{\figure{\begin{tikzpicture}[overlay,x=1em,y=1em]\draw[line width=0.35pt] (-0.10,-0.16) -- (0,-0.16);\end{tikzpicture}}}}\uline{ien} et travaillant avec la $(n+1)$-catégorie (sans \unsure{structures|structure} \unsure{supplémentaires|supplémentaire} dessus!) des $n$-champs loc. constants\raisebox{0pt}[0pt][0pt]{\parbox{0pt}{\figure{\begin{tikzpicture}[overlay,x=1em,y=1em]\fill (0.08,0.62) circle[radius=0.04em];\end{tikzpicture}}}}
% L23: preferred live raised addition, with curved attachment before pour; cancellation remains a state alternative.
\inserted{\unsure{sur $X$}} pour tout $n$
